% frontmatter.tex -- title block and abstract, transcribed from the Typst source.
%
% Changelog (reverse-chronological)
%   2026-09-01  Created.
%
% tools/typ2tex.sh strips everything before the first "= " heading (the Typst
% preamble imports theorion, which pandoc's evaluator cannot run), and the title
% block, author block, ORCID, affiliation, e-mail, abstract and title footnote
% all live in that stripped region.  They are transcribed here BY HAND from
% article/ring-of-k-commuting-2x2-matrices.typ, which makes this file a second
% copy that can drift from the canonical source -- it did, within a day of being
% written.  Edit the .typ first, then mirror the change here, and rely on
% make-arxiv.sh's internal-matter check to catch what mirroring misses.  The
% Acknowledgements paragraph is NOT here: it sits after the bibliography in the
% Typst and survives the conversion on its own.

\definecolor{orcidgreen}{HTML}{A6CE39}
\newcommand{\orcidid}[1]{%
  \href{https://orcid.org/#1}{%
    \textcolor{orcidgreen}{\fontsize{8}{9}\selectfont iD}%
    \,\fontsize{8}{9}\selectfont #1}}

% NB: the title is NOT wrapped in a minipage.  A \footnote inside a minipage is
% typeset at the foot of the minipage, not the page -- the title footnote would
% land immediately under the title.
\begingroup
\centering
% "plus 1fil" on BOTH sides, not bare lengths: \centering works by setting
% \leftskip/\rightskip to 0pt plus 1fil, so overwriting them with rigid lengths
% narrows the measure but also destroys the centring -- the title's last line
% came out stretched across the full measure.
\leftskip=0.05\textwidth plus 1fil \rightskip=0.05\textwidth plus 1fil
    \fontsize{15}{18}\selectfont\bfseries
    The coordinate ring of the $k$-fold iterated commutator locus for
    $2 \times 2$ matrices%
    \par
\endgroup

\begingroup
\centering
\vspace{0.8em}
{\fontsize{11}{13}\selectfont Jan Snellman\hspace{0.25em}%
  \orcidid{0009-0002-6676-5068}\par}

\vspace{0.35em}

{\fontsize{9}{11}\selectfont Matematiska Institutionen, Linköpings
  Universitet, 581 83 Linköping, Sweden\par}

\vspace{0.15em}

{\fontsize{9}{11}\selectfont
  \href{mailto:jan.snellman@liu.se}{jan.snellman@liu.se}\par}
\endgroup

\vspace{1.0em}

\begin{center}
\begin{minipage}{\dimexpr\textwidth-3.2cm\relax}
  \fontsize{9.2}{11.5}\selectfont
  \begin{center}\textbf{Abstract}\end{center}
  \vspace{0.25em}
  For $2 \times 2$ matrices $A_1,\dots,A_k$, write $[A_1,\dots,A_k]$ for the
  left-normed iterated commutator $[\dots[[A_1,A_2],A_3],\dots,A_k]$, and $I_k$
  for the ideal, in the $3k$-variable reduced-coordinate polynomial ring $R_k$,
  cutting out its vanishing locus. We prove, for every $k \geq 2$ over any field
  of characteristic $\neq 2$, and as four independently-established results
  rather than one bundled claim: $I_k$ has codimension 2; $I_k$ has exactly 3
  minimal generators; $R_k / I_k$ is Cohen-Macaulay; and $I_k$ is radical. The
  last of these, together with an explicit component count resting on a
  non-containment argument, assembles into the \textbf{Primary Decomposition
  Theorem}: $I_k = P_2 \cap \cdots \cap P_k$ is an irredundant primary
  decomposition into exactly $k-1$ primes, following an explicit recursive
  block-involvement pattern. The proof identifies $I_k$ as the ideal of
  $2 \times 2$ minors of an explicit $2 \times 3$ matrix (a determinantal ideal,
  not merely determinantal-\textit{looking}), and invokes classical
  determinantal-ideal theory (Bruns-Vetter) and an explicit rank-2 Jacobian
  witness on every component (Serre's criterion) for the algebraic and
  radicality halves respectively.
\end{minipage}
\end{center}

\vspace{0.8em}
